\chapter{الكحولات: تنشيط مجموعة OH}\label{ch:b1:alcohol-activation}

الكحولات في كل مكان: تُصنع بالأطنان من الألكينات ومن
التخمير، وتعطيها كل إضافة \omterm{def:b1:organomagnesium:organometallic}{لكاشف غرينيار} في الفصل
السابق. ومع ذلك لا يخضع الكحول للاستبدال كما يخضع الهالوجينو ألكان:
إذ كان على مجموعة OH أن تغادر على هيئة أيون الهيدروكسيد، وهو \omterm{def:b1:predominant-reaction:strong-acid}{قاعدة قوية}
\omterm{def:b1:electronic-effects:nucleophile}{ومجموعة مغادرة} رديئة جدًا (\cref{ch:b1:nucleophilic-substitution}).
وللكيميائيين ثلاثة أجوبة. انزع بروتون OH بدلًا من ذلك، فيصير
\omterm{def:b1:alcohol-activation:alkoxide}{الألكوكسيد} \omterm{def:b1:electronic-effects:nucleophile}{محبًا للنواة} قويًا. أو أضف بروتونًا، فتصير \omterm{def:b1:electronic-effects:nucleophile}{المجموعة المغادرة}
الماء. أو استبدل بالهيدروجين مجموعة تجعل
الأكسجين جزءًا من \omterm{def:b1:electronic-effects:nucleophile}{مجموعة مغادرة} ممتازة، دون أن تمس
الكربون. يتتبع هذا الفصل الطرق الثلاث، وما يرافقها من خيارات في
الآلية والكيمياء الفراغية.

\begin{recall}
SN1 وSN2 وكيميائهما الفراغية (\cref{ch:b1:nucleophilic-substitution})؛
وE1 وE2 \omterm{prop:b1:elimination:zaitsev}{وقاعدة زايتسيف} (\cref{ch:b1:elimination})؛ وثوابت الحمضية
وطريقة التفاعل السائد (\cref{ch:b1:predominant-reaction})؛
وأثر مجموعات الألكيل في حمضية الكحولات
(\cref{exo:b1:electronic-effects:11}).
\end{recall}

\section{الكحولات أحماضًا: الألكوكسيدات}

\begin{definition}[الألكوكسيد]\label{def:b1:alcohol-activation:alkoxide}
\emph{الألكوكسيد}\index{ألكوكسيد} هو القاعدة المرافقة \ce{RO-} لكحول
\ce{ROH}: الميثوكسيد \ce{CH3O-}، والإيثوكسيد \ce{CH3CH2O-}،
وثالثي البوتوكسيد \ce{(CH3)3CO-}. والألكوكسيدات \omterm{def:b1:predominant-reaction:strong-acid}{قواعد قوية} \omterm{def:b1:electronic-effects:nucleophile}{ومحبات
للنواة} جيدة.
\end{definition}

\begin{proposition}[حمضية الكحولات]\label{prop:b1:alcohol-activation:alcohol-acidity}
الكحولات أحماض أضعف من الماء أو تقاربه ضعفًا: $\mathrm{p}K_a$ يساوي
15.5 (الميثانول)، و15.9 (الإيثانول)، و17.1 (البروبان-2-أول)، و19.2
(2-ميثيل البروبان-2-أول)، مقابل 14.0 للمزدوجة $\ce{H2O}/\ce{OH-}$.
ولذلك لا يزيل الهيدروكسيد بروتوناتها إلا جزئيًا؛ والتحويل التام
إلى \omterm{def:b1:alcohol-activation:alkoxide}{الألكوكسيد} يحتاج إلى قاعدة أقوى أو إلى فلز قلوي.
\end{proposition}
% ledger: pka:methanol, pka:ethanol, pka:2-propanol, pka:tBuOH

\begin{proof}
للتفاعل \ce{ROH + OH- <=> RO- + H2O} الثابت $K = 10^{14.0 -
\mathrm{p}K_a(\ce{ROH})}$، بين $10^{-1.5}$ و$10^{-5}$: فهو غير ملائم.
أما فلز الصوديوم فيختزل البروتون، \ce{2ROH + 2Na -> 2RO- + 2Na+ + H2}،
وأيون الهيدريد في هيدريد الصوديوم يأخذه أخذًا لا عكوسًا،
\ce{ROH + H- -> RO- + H2}: فيفلت الهيدروجين، ويكون التفاعل
تامًا مهما كان ثابت توازنه.
\end{proof}

\begin{inthelab}[هيدريد الصوديوم]
يُباع هيدريد الصوديوم مسحوقًا رماديًا مشتتًا في زيت معدني، يحميه
من رطوبة الهواء. ويتفاعل بعنف مع الماء،
فيطلق هيدروجينًا قابلًا للاشتعال؛ ولذلك يوزن بسرعة، ويُضاف على دفعات
إلى الكحول الجاف تحت غاز خامل، ويدل فوران الهيدروجين على
تقدم نزع البروتون. ويُهدم الفائض من الهيدريد في
النهاية بإضافة بطيئة لكحول، لا للماء أبدًا.
\end{inthelab}

\section{تخليق ويليامسون للإيثرات}

\begin{definition}[تخليق ويليامسون للإيثرات]\label{def:b1:alcohol-activation:williamson}
\emph{تخليق ويليامسون للإيثرات}\index{تخليق ويليامسون للإيثرات}
يحضّر إيثرًا $\mathrm{R{-}O{-}R'}$ بتفاعل SN2 \omterm{def:b1:alcohol-activation:alkoxide}{لألكوكسيد} \ce{RO-}
على هالوجينو ألكان (أو \omterm{def:b1:alcohol-activation:sulfonate}{إستر سلفونات}) $\mathrm{R'{-}X}$.
\end{definition}

\begin{omfigure}
\setchemfig{atom sep=2.1em}
\schemestart
\chemfig{CH_3CH_2-@{o}O^{-}}\,\chemfig{Na^{+}}
\+
\chemfig{@{c}CH_3-[@{ci}]@{i}I}
\arrow{->}
\chemfig{CH_3CH_2-O-CH_3}
\+
\chemfig{Na^{+}}\,\chemfig{I^{-}}
\schemestop

\medskip
\chemmove{\draw[->, thick, shorten <=2pt, shorten >=2pt](o) .. controls +(80:8mm) and +(100:8mm) .. (c);
          \draw[->, thick, shorten <=2pt, shorten >=2pt](ci) .. controls +(90:5mm) and +(90:5mm) .. (i);}

{\small تخليق ويليامسون للميثوكسي إيثان: يهاجم الإيثوكسيد
كربون اليودوميثان من الجهة المقابلة لليود (SN2).}
\end{omfigure}

\begin{method}[اختيار الشريكين]\label{met:b1:alcohol-activation:williamson-choice}
يمكن قطع الإيثر $\mathrm{R{-}O{-}R'}$ على أيٍّ من جانبي الأكسجين.
\begin{enumerate}
  \item اكتب الزوجين الممكنين: \ce{RO-} مع $\mathrm{R'X}$،
        و$\mathrm{R'O^-}$ مع \ce{RX}.
  \item احتفظ بالزوج الذي يكون فيه الهاليد ميثيليًا أو أوليًا: \omterm{def:b1:nucleophilic-substitution:sn2}{فالآلية SN2} تحتاج إلى
        كربون غير معاق.
  \item ارفض الزوج ذا الهاليد الثالثي: \omterm{def:b1:alcohol-activation:alkoxide}{فالألكوكسيد}، وهو \omterm{def:b1:predominant-reaction:strong-acid}{قاعدة قوية}،
        سيُحدث حذفًا (E2) بدلًا من ذلك.
  \item ومع هاليد ثانوي، توقّع خليطًا من الإيثر والألكين.
\end{enumerate}
\end{method}

\begin{example}[إيثر غير متناظر]\label{ex:b1:alcohol-activation:williamson}
2-ميثوكسي-2-ميثيل البروبان (ميثيل ثالثي بوتيل الإيثر): يعطيه ثالثي بوتوكسيد الصوديوم
مع اليودوميثان بنقاء؛ أما ميثوكسيد الصوديوم مع البروميد
الثالثي، 2-برومو-2-ميثيل البروبان، فيعطي 2-ميثيل البروبين \omterm{def:b1:elimination:elimination}{بالآلية E2}. الضخامة في جانب
\omterm{def:b1:alcohol-activation:alkoxide}{الألكوكسيد}، والصغر في جانب الهاليد.
\end{example}

\section{تنشيط مجموعة OH}

\begin{definition}[إسترات السلفونات]\label{def:b1:alcohol-activation:sulfonate}
يتكوّن \emph{إستر السلفونات}\index{إستر السلفونات} $\mathrm{R{-}O{-}SO_2{-}R'}$
من كحول وكلوريد سلفونيل $\mathrm{R'SO_2Cl}$ بوجود
قاعدة (البيريدين). و\emph{التوزيلات}\index{توزيلات}، \ce{R-OTs}، تأتي
من كلوريد 4-ميثيل بنزين السلفونيل (كلوريد التوزيل)؛
و\emph{الميزيلات}\index{ميزيلات}، \ce{R-OMs}، من كلوريد ميثان السلفونيل.
وأنيون السلفونات $\mathrm{R'SO_3^-}$، القاعدة المرافقة \omterm{def:b1:predominant-reaction:strong-acid}{لحمض قوي}
والمثبَّت بالرنين على ثلاث ذرات أكسجين، \omterm{def:b1:electronic-effects:nucleophile}{مجموعة مغادرة} ممتازة.
\end{definition}

\begin{omfigure}
\setchemfig{atom sep=2em}
\schemestart
\chemfig{R-O-H}
\+
\chemfig{Cl-S(=[2]O)(=[6]O)-[:0]*6(-=-(-CH_3)=-=)}
\arrow{->[بيريدين]}
\chemfig{R-O-S(=[2]O)(=[6]O)-[:0]*6(-=-(-CH_3)=-=)}
\schemestop
\par\vspace{6pt}

{\small توزلة كحول. تحل الرابطة O--S محل الرابطة O--H؛ ولا تُمس
الرابطة C--O في الكحول. ويأخذ البيريدين HCl المتكوّن.}
\end{omfigure}

\begin{proposition}[الترتيب الفراغي عبر التوزلة والاستبدال]\label{prop:b1:alcohol-activation:tosylation-retention}
توزلة كحول على كربون فراغي المركز تحفظ ترتيبه الفراغي
(احتفاظ)؛ والاستبدال SN2 اللاحق \omterm{def:b1:alcohol-activation:sulfonate}{للتوزيلات} يقلبه.
والخطوتان معًا تستبدلان بمجموعة OH \omterm{def:b1:electronic-effects:nucleophile}{محبًا للنواة} مع انقلاب.
\end{proposition}

\begin{proof}
في التوزلة، الروابط التي تتغير هي O--H وS--Cl: فأكسجين
الكحول يهاجم الكبريت ويذهب هيدروجينه إلى القاعدة. ولا
تنكسر أي رابطة مع الكربون الفراغي المركز، فيبقى الترتيب حوله
نفسه. وفي خطوة SN2، تنكسر الرابطة C--O بينما يرتبط \omterm{def:b1:electronic-effects:nucleophile}{المحب للنواة}
من الجهة المقابلة (\cref{prop:b1:nucleophilic-substitution:sn2-inversion}).
\end{proof}

\begin{omfigure}
\begin{tikzpicture}[font=\small, >={Stealth}]
  \node (a) at (0,0) {\chemfig{C(-[:150]H_3C)(<[:210]H)(<:[:250]C_2H_5)-OH}};
  \node at (0,-1.4) {\foreignlanguage{arabic}{بوتان-\babelsublr{2}-أول-\babelsublr{\cip{S}}}};
  \draw[->, thick] (1.3,0) -- (2.5,0) node[midway, above] {TsCl} node[midway, below, font=\scriptsize] {احتفاظ};
  \node (b) at (3.9,0) {\chemfig{C(-[:150]H_3C)(<[:210]H)(<:[:250]C_2H_5)-OTs}};
  \node at (3.9,-1.4) {\foreignlanguage{arabic}{توزيلات \babelsublr{\cip{S}}}};
  \draw[->, thick] (5.3,0) -- (6.6,0) node[midway, above] {\ce{CN-}} node[midway, below, font=\scriptsize] {\foreignlanguage{arabic}{\babelsublr{SN2}، انقلاب}};
  \node (c) at (8.3,0) {\chemfig{[:0]NC-C(-[:30]CH_3)(<[:-30]H)(<:[:-70]C_2H_5)}};
  \node at (8.3,-1.4) {\foreignlanguage{arabic}{نتريل \babelsublr{\cip{R}}}};
\end{tikzpicture}

{\small من بوتان-2-أول-\cip{S} إلى 2-ميثيل بوتان نتريل: تحفظ التوزلة
\omterm{def:b1:stereochemistry-in-depth:isomers}{الترتيب الفراغي}، وتقلبه خطوة SN2 مع السيانيد. وإجمالًا:
حلّت CN محل OH مع انقلاب.}
\end{omfigure}

ويمكن بلوغ الهدف نفسه في وسط حمضي: فمجموعة OH، بعد برتنتها، تغادر على هيئة
ماء، \ce{R-OH2+ -> R+ + H2O} (SN1، E1) أو تحت الهجوم (SN2). غير أن الحمض يقصر
\omterm{def:b1:electronic-effects:nucleophile}{المحبات للنواة} على تلك التي تصمد فيه --- أيونات الهاليد، والماء،
والكحولات --- ويفتح الباب \omterm{def:b1:electronic-effects:intermediates}{للكاتيونات الكربونية} وتفاعلاتها الجانبية.

\section{التحويل إلى هالوجينو ألكانات}

\begin{proposition}[الكحولات وهاليدات الهيدروجين]\label{prop:b1:alcohol-activation:hx-by-class}
تعطي الكحولات الثالثية الهالوجينو ألكان مع حمض الهيدروكلوريك المركّز
عند درجة حرارة الغرفة (SN1)؛ وتحتاج الكحولات الأولية إلى حمض الهيدروبروميك أو
حمض الهيدرويوديك وإلى التسخين (SN2 على الكحول المتبرتن)؛ وتتفاعل الكحولات
الثانوية بسرعات وسطى، وكثيرًا ما يكون ذلك بالآليتين.
\end{proposition}

\begin{proof}
تحوّل البرتنة مجموعة OH إلى \ce{-OH2+}. وفي الكحول الثالثي يغادر الماء
بسهولة، فيتكوّن \omterm{def:b1:electronic-effects:intermediates}{كاتيون كربوني} ثابت يأسره الهاليد (SN1): فحتى
\omterm{def:b1:electronic-effects:nucleophile}{المحب للنواة} المتواضع \ce{Cl-} يكفي. أما \omterm{def:b1:electronic-effects:intermediates}{الكاتيون الكربوني} الأولي فغير ثابت
إلى حد كبير؛ فيجب أن يدفع الهاليد الماء إلى الخارج من الخلف (SN2)، وهذا يحتاج إلى
\omterm{def:b1:electronic-effects:nucleophile}{محب للنواة} جيد (\ce{Br-}، \ce{I-}، أفضل في الأوساط البروتونية من \ce{Cl-}،
\cref{prop:b1:electronic-effects:nucleophilicity-basicity}) وإلى الحرارة.
\end{proof}

\begin{example}[اختبار للتصنيف]\label{ex:b1:alcohol-activation:lucas}
خليط من حمض الهيدروكلوريك المركّز وكلوريد الزنك (\omterm{def:b1:lewis-resonance-vsepr:lewis-acid}{حمض لويس}
يساعد مجموعة OH على المغادرة) يعكّر الكحول الثالثي فورًا ---
إذ ينفصل الكلوريد غير الذواب --- والكحول الثانوي في غضون دقائق،
ولا يكاد يعكّر الكحول الأولي عند درجة حرارة الغرفة: إنه ترتيب
ثبات \omterm{def:b1:electronic-effects:intermediates}{الكاتيونات الكربونية} وقد صار مرئيًا في أنبوب اختبار.
\end{example}

\section{نزع الماء وتكوين الإيثر في الوسط الحمضي}

\begin{definition}[نزع الماء من كحول]\label{def:b1:alcohol-activation:dehydration}
\emph{نزع الماء}\index{نزع الماء} من كحول هو
حذف الماء، بحفز \omterm{def:b1:predominant-reaction:strong-acid}{حمض قوي} (حمض الكبريتيك أو حمض الفوسفوريك)
وبالحرارة، فيتكوّن ألكين: $\mathrm{R_2CH{-}CR'_2OH} \to \mathrm{R_2C{=}CR'_2} + \ce{H2O}$.
\end{definition}

\begin{omfigure}
\setchemfig{atom sep=2.1em}
\schemestart
\chemfig{C(-[2]CH_3)(-[:210]H_3C)(-[:-30]CH_2CH_3)-OH}
\arrow{->[\ce{H+}]}
\chemfig{C(-[2]CH_3)(-[:210]H_3C)(-[:-30]CH_2CH_3)-OH_2^{+}}
\arrow{->[\ce{-H2O}]}
\chemfig{C^{+}(-[2]CH_3)(-[:210]H_3C)(-[:-30]CH_2CH_3)}
\schemestop

\medskip
\schemestart
\chemfig{C^{+}(-[2]CH_3)(-[:210]H_3C)(-[:-30]CH_2CH_3)}
\arrow{->[\ce{-H+}][رئيسي]}
\chemfig{C(-[2]CH_3)(-[:210]H_3C)=[:-30]CHCH_3}
\schemestop

\medskip
{\small \omterm{def:b1:alcohol-activation:dehydration}{نزع الماء} من 2-ميثيل بوتان-2-أول بحفز حمضي (E1): برتنة،
ثم فقد الماء نحو \omterm{def:b1:electronic-effects:intermediates}{الكاتيون الكربوني} الثالثي، ثم فقد بروتون. ونزع
البروتون من \ce{CH2} يعطي 2-ميثيل بوت-2-ين الثلاثي الاستبدال
(زايتسيف، الرئيسي)؛ ونزعه من مجموعة ميثيل يعطي 2-ميثيل بوت-1-ين (الثانوي).}
\end{omfigure}

\begin{proposition}[إيثر أم ألكين]\label{prop:b1:alcohol-activation:ether-vs-alkene}
يعطي الكحول الأولي المسخَّن مع حمض الكبريتيك أساسًا إيثرًا متناظرًا
عند درجة حرارة معتدلة (جزيء يستبدل جزيئًا آخر \omterm{def:b1:nucleophilic-substitution:sn2}{بالآلية SN2})
وأساسًا ألكينًا عند درجة حرارة أعلى؛ وتعطي الكحولات الثانوية والثالثية
ألكينات.
\end{proposition}

\begin{proof}
في الكحول الأولي، يمكن أن يُهاجَم الكحول المتبرتن إما عند
الكربون بجزيء كحول ثانٍ (SN2، فيعطي \ce{R-O-R} بعد فقد
بروتون) وإما عند هيدروجين $\beta$ (E2). والحذف يكوّن جزيئات أكثر
مما يستهلك، فهو محابى كلما ارتفعت درجة الحرارة
(\cref{prop:b1:elimination:sn-vs-e})؛ ويسود الاستبدال حين تكون
أدنى. أما الكحولات الثانوية والثالثية فتتأين إلى \omterm{def:b1:electronic-effects:intermediates}{كاتيونات كربونية}، تفقد
بروتونًا بسهولة أكبر مما يأسرها \omterm{def:b1:electronic-effects:nucleophile}{محب للنواة} ضعيف، ولا سيما
عند التسخين؛ ويُزال الماء المتكوّن أيضًا بتقطير الألكين
المتطاير.
\end{proof}

\section{تمارين}

\begin{exercise}[$\star$]\label{exo:b1:alcohol-activation:1}
اكتب تفاعلات الإيثانول مع الصوديوم، ومع هيدريد الصوديوم، ومع
هيدروكسيد الصوديوم؛ واحسب ثابت التفاعل الأخير ($\mathrm{p}K_a$
للإيثانول يساوي 15.9).
\end{exercise}

\begin{exercise}[$\star$]\label{exo:b1:alcohol-activation:2}
أعطِ نواتج: ميثوكسيد الصوديوم مع 1-بروموبروبان؛ وإيثوكسيد
الصوديوم مع اليودوميثان؛ وفينوكسيد الصوديوم مع البروموإيثان.
\end{exercise}

\begin{exercise}[$\star$]\label{exo:b1:alcohol-activation:3}
اكتب توزلة البروبان-1-أول وتفاعل \omterm{def:b1:alcohol-activation:sulfonate}{التوزيلات} مع
يوديد الصوديوم. وما الميزة مقارنةً بالتفاعل المباشر
للكحول مع يوديد الصوديوم؟
\end{exercise}

\begin{exercise}[$\star$]\label{exo:b1:alcohol-activation:4}
أعطِ الألكينات المتكوّنة \omterm{def:b1:alcohol-activation:dehydration}{بنزع الماء} الحمضي من البوتان-2-أول ومن
2-ميثيل البنتان-2-أول، والألكين الرئيسي.
\end{exercise}

\begin{exercise}[$\star\star$]\label{exo:b1:alcohol-activation:5}
اقترح تخليق ويليامسون للمركب 1-إيثوكسي-2-ميثيل البروبان وعلّل
اختيارك للشريكين. ولماذا كان الاختيار الآخر سيعطي ألكينًا؟
\end{exercise}

\begin{exercise}[$\star\star$]\label{exo:b1:alcohol-activation:6}
يُحوَّل الأوكتان-2-أول-\cip{R} إلى توزيلاته، ثم يُعالج بإيثوكسيد
الصوديوم في الإيثانول. أعطِ \omterm{def:b1:stereochemistry-in-depth:isomers}{الترتيب الفراغي} للإيثر. وماذا
كان سيُحصل عليه بتسخين الكحول نفسه مع حمض الكبريتيك؟
\end{exercise}

\begin{exercise}[$\star\star$]\label{exo:b1:alcohol-activation:7}
اكتب \omterm{def:b1:elementary-steps:elementary-step}{آلية تفاعل} 2-ميثيل البروبان-2-أول مع
حمض الهيدروكلوريك المركّز وفسّر لماذا كان التفاعل سريعًا عند
درجة حرارة الغرفة.
\end{exercise}

\begin{exercise}[$\star\star$]\label{exo:b1:alcohol-activation:8}
اكتب آلية تكوّن الإيثوكسي إيثان من الإيثانول في
حمض الكبريتيك، والحذف المنافس.
\end{exercise}

\begin{exercise}[$\star\star$]\label{exo:b1:alcohol-activation:9}
في اختبار التصنيف في هذا الفصل، أي كحول من بين
البوتان-1-أول والبوتان-2-أول و2-ميثيل البروبان-2-أول يتعكّر أولًا؟ اكتب
الناتج.
\end{exercise}

\begin{exercise}[$\star\star\star$]\label{exo:b1:alcohol-activation:10}
يعطي 3,3-ثنائي ميثيل البوتان-2-أول، عند تسخينه في وسط حمضي، أساسًا
2,3-ثنائي ميثيل بوت-2-ين، الذي يختلف هيكله الكربوني عن هيكل
الكحول. اقترح آلية فيها انتقال 1,2 لمجموعة ميثيل في
\omterm{def:b1:electronic-effects:intermediates}{الكاتيون الكربوني}، وعلّل هذا الانتقال.
\end{exercise}

\begin{exercise}[$\star\star\star$]\label{exo:b1:alcohol-activation:11}
صمّم تخليقًا للمركب 2-ميثوكسي بوتان-\cip{S} انطلاقًا من بوتان-2-أول-\cip{R}،
وآخر انطلاقًا من بوتان-2-أول-\cip{S}، يبقى \omterm{def:b1:stereochemistry-in-depth:isomers}{الترتيب الفراغي} في كل منهما تحت
السيطرة.
\end{exercise}

\begin{exercise}[$\star\star\star$]\label{exo:b1:alcohol-activation:12}
بيّن أن التوازن \ce{2CH3CH2OH <=> (CH3CH2)2O + H2O} يمكن
دفعه نحو الإيثر بإزالة الماء باستمرار، مستعملًا حاصل
التفاعل (\cref{ch:b1:extent-q-and-k}).
\end{exercise}

\section{مسألة: طريقان إلى إيثر}

\begin{problem}[{مسألة نهاية الأسبوع --- طريقا ويليامسون إلى ميثيل
ثالثي بوتيل الإيثر، وإيثر مضبوط الكيمياء الفراغية من كحول نشط
ضوئيًا، ونزع الماء من كحول ثالثي، وكتلة
الإيثر المحصَّل عليها}]
\label{pb:b1:alcohol-activation:1}
أُنتج 2-ميثوكسي-2-ميثيل البروبان (ميثيل ثالثي بوتيل الإيثر، MTBE) على
نطاق واسع مادةً مضافة إلى الوقود. الكتل المولية (\unit{g/mol}): \ce{C4H10O}
74.0، \ce{C5H12O} 88.0، \ce{CH3I} 141.9.

\medskip
\noindent\textbf{الجزء الأول --- MTBE بتخليق ويليامسون.}
\begin{enumerate}
  \item اكتب زوجي الشريكين (\omterm{def:b1:alcohol-activation:alkoxide}{ألكوكسيد} وهاليد) اللذين يمكن أن
        يعطيا MTBE.
  \item أي الزوجين يفشل؟ اكتب التفاعل الذي يحدث بدلًا من ذلك.
  \item اكتب التفاعل الناجح وآليته.
  \item كيف يُحضَّر ثالثي بوتوكسيد الصوديوم من 2-ميثيل البروبان-2-أول؟
        ولماذا لا يُحضَّر بهيدروكسيد الصوديوم؟
  \item لماذا يُجرى التفاعل في الكحول الجاف أو في \omterm{def:b1:intermolecular-forces-solvents:solvent-classes}{مذيب لابروتوني}؟
  \item هل الناتج \omterm{def:b1:stereochemistry-in-depth:stereocentre}{كيرالي}؟
\end{enumerate}

\medskip
\noindent\textbf{الجزء الثاني --- إيثر نشط ضوئيًا.}
\begin{enumerate}[resume]
  \item يُحوَّل البوتان-2-أول-\cip{R} إلى توزيلاته. اكتب
        التفاعل وأعطِ \omterm{def:b1:stereochemistry-in-depth:isomers}{الترتيب الفراغي} \omterm{def:b1:alcohol-activation:sulfonate}{للتوزيلات}.
  \item تتفاعل \omterm{def:b1:alcohol-activation:sulfonate}{التوزيلات} مع ميثوكسيد الصوديوم. اكتب الناتج
        وترتيبه الفراغي.
  \item أي \omterm{def:b1:stereochemistry-in-depth:isomers}{ترتيب فراغي} للكحول يعطي 2-ميثوكسي بوتان-\cip{R}
        بالطريق نفسه؟
  \item طريق آخر: \omterm{def:b1:alcohol-activation:alkoxide}{ألكوكسيد} البوتان-2-أول-\cip{R} مع
        اليودوميثان. ما \omterm{def:b1:stereochemistry-in-depth:isomers}{الترتيب الفراغي} للناتج؟ ولماذا؟
  \item أي تفاعل جانبي ينافس في السؤال 8، ولماذا كان
        محدودًا هنا؟
  \item لماذا كان تسخين البوتان-2-أول-\cip{R} مع الميثانول في حمض
        الكبريتيك سيعطي إيثرًا شبه راسيمي، إن أعطى إيثرًا أصلًا؟
\end{enumerate}

\medskip
\noindent\textbf{الجزء الثالث --- \omterm{def:b1:alcohol-activation:dehydration}{نزع الماء} من كحول ثالثي.}
\begin{enumerate}[resume]
  \item اكتب آلية \omterm{def:b1:alcohol-activation:dehydration}{نزع الماء} الحمضي من
        2-ميثيل بوتان-2-أول.
  \item أعطِ الألكينين وتوقّع الرئيسي منهما.
  \item لماذا لا يتكوّن إيثر من هذا الكحول؟
  \item لماذا يُقطَّر 2-ميثيل بوت-2-ين إلى الخارج كلما تكوّن؟
  \item ارسم منحنى الطاقة لخطوة E1 وعلّم
        \omterm{def:b1:elementary-steps:rds}{الخطوة المحددة للسرعة}.
  \item هل يمكن \omterm{def:b1:alcohol-activation:dehydration}{نزع الماء} من كحول أولي بالآلية نفسها؟
\end{enumerate}

\medskip
\noindent\textbf{الجزء الرابع --- الكتل.}
\begin{enumerate}[resume]
  \item يُحوَّل \qty{10.0}{g} من 2-ميثيل البروبان-2-أول إلى
        \omterm{def:b1:alcohol-activation:alkoxide}{الألكوكسيد}. احسب كمية المادة.
  \item ما كتلة اليودوميثان اللازمة لفائض قدره \qty{10}{\percent}؟
  \item احسب الكتلة النظرية للمركب MTBE.
  \item احسب الكتلة المحصَّل عليها \omterm{def:b1:extent-q-and-k:fractional-extent}{بمردود} \qty{75}{\percent}.
\end{enumerate}
\end{problem}
